\section{Evaluating the decision layer}
\label{sec:toolselection}

The preceding sections measure the physical chain and the repair loop. One link
remained unevaluated, yet central: faced with a request written in natural
language, does the supervisor choose the right tool?

\subsection{Protocol}

Twenty-six requests were written, covering the six tools exposed to the agent,
plus four requests for which \emph{no} tool should be called — a greeting, a
general-knowledge question, a note of thanks. Triggering a tool on these is a
false positive just as much as a wrong choice is.

The tools are not executed. They are bound to the model, and it is the
\emph{intent} produced — the list of requested calls — that is recorded. This
protocol measures selection exactly, with no side effects: no case is created,
no solver is launched. It takes one pass per request.

It should be stressed that this evaluation was not practicable before the
hardware acceleration described in section~\ref{sec:latence}. At four minutes
per invocation, these twenty-six requests would have taken almost two hours;
they now run in under a minute, which allows several iterations within a single
working session.

\subsection{Results and diagnosis}

The first campaign gives an overall accuracy of \SI{77}{\percent}. The
breakdown by tool, however, is very uneven (table~\ref{tab:toolsel}) and points
to a single culprit: the safety analyser is correctly selected in only one case
out of four.

\begin{table}[htbp]
\centering
\caption[Tool-selection accuracy]{Tool-selection accuracy, before and after
rewording the descriptions. A case counts as correct if the expected tool is
called first, or if no tool is called when none is expected.}
\label{tab:toolsel}
\begin{tabularx}{\textwidth}{@{}Xccc@{}}
\toprule
\textbf{Expected tool} & \textbf{$n$} & \textbf{Before} & \textbf{After}\\
\midrule
Case creation & 4 & 4/4 & 4/4\\
Solver execution & 4 & 4/4 & 4/4\\
Post-processing & 4 & 3/4 & 3/4\\
Safety analysis & 4 & \textcolor{cAlert}{\textbf{1/4}} & \textbf{4/4}\\
Surrogate & 3 & 2/3 & 2/3\\
Documentation & 3 & 3/3 & 3/3\\
No tool expected & 4 & 3/4 & 3/4\\
\midrule
\textbf{Total} & \textbf{26} & \textbf{20/26 (\SI{77}{\percent})} &
  \textbf{23/26 (\SI{88}{\percent})}\\
\bottomrule
\end{tabularx}
\end{table}

Examining the three failures is enlightening. On ``check whether a safety
criterion is exceeded'', the agent calls the executor; on ``analyse the PCMI
margins and tell me when the limit will be reached'', it calls the surrogate; on
``is there a risk of fuel melting?'', it calls nothing.

The cause is not the model, but the descriptions supplied. Those of the safety
analyser and of the surrogate promised exactly the same thing — safety margins,
per-criterion status, PCMI, gap closure — without ever stating what really
distinguishes them. Yet that criterion is simple and unambiguous: \textbf{the
analyser requires an already-simulated case on disk, whereas the surrogate takes
only numerical parameters}.

\subsection{Correction and re-measurement}

The two descriptions were rewritten to state this condition of use explicitly,
and to name each other as alternatives not to be confused. No other change was
made — neither to the model, nor to the tool code, nor to the requests.

The accuracy of the safety analyser rises from \SI{25}{\percent} to
\SI{100}{\percent}, and the overall accuracy from \SI{77}{\percent} to
\SI{88}{\percent}.

The lesson goes beyond this particular case. In a tool-based architecture, a
tool description is not documentation: it is the \emph{programming interface}
through which the model decides. Two tools whose descriptions overlap are, from
the model's point of view, indistinguishable — and no amount of model quality
compensates for that ambiguity. It is also good news in terms of cost: eleven
accuracy points were gained by rewriting two paragraphs, with no retraining and
no change of model.

The three remaining errors do not stem from this cause. Two concern wordings
where the boundary is genuinely debatable — ``quickly estimate the maximum
temperature'' may legitimately call either the surrogate or post-processing —
and the third is a documentation call on a general-knowledge question, a benign
false positive.

\subsection{From selection to execution: a blind spot}
\label{sec:angle-mort}

The evaluation above measures which tool the agent \emph{chooses}. It says
nothing about what happens when that call is actually \emph{executed}, since the
protocol reads the intent without ever triggering it. That choice, justified in
order to isolate the decision, turned out to be a blind spot: the first attempt
at end-to-end execution failed completely, on a defect invisible to the
selection benchmark.

\textbf{A tool schema incompatible with actual use.} The creation tool declared
its \texttt{params} parameter as a \emph{string} containing JSON. Yet a language
model spontaneously produces an \emph{object} for an argument so named.
Validation rejected it, the agent received an error message, and replayed the
identical call. Measured: \num{30} tool calls, \num{53} messages,
\SI{179}{\second}, \textbf{no case created at all}. Accepting both forms removed
the loop.

\textbf{A context window that was too narrow.} A second defect manifested itself
as freezes on multi-step requests, while a simple question was answered in one
second. Measurement showed that the system prompt and the tool schemas alone
consume \textbf{$\sim$\num{1900} tokens, that is \SI{47}{\percent} of the
\num{4096}-token window} used by default — before any exchange. A single tool
result was then enough to overflow it. The window was raised to \num{16384}
tokens.

Table~\ref{tab:bout-en-bout} gives the combined effect of these two corrections
on the same request.

\begin{table}[htbp]
\centering
\caption[End-to-end execution]{End-to-end execution of a four-step request
(create, simulate, post-process, analyse safety), before and after correcting
the tool schema and the context window.}
\label{tab:bout-en-bout}
\begin{tabularx}{0.92\textwidth}{@{}Xcc@{}}
\toprule
\textbf{Indicator} & \textbf{Before} & \textbf{After}\\
\midrule
Tool calls & \num{30} & \textbf{\num{5}}\\
Messages exchanged & \num{53} & \textbf{\num{11}}\\
Total duration & \SI{179}{\second} & \textbf{\SI{39}{\second}}\\
Case created and simulated & \textcolor{cAlert}{\textbf{no}} & \textbf{yes}\\
\bottomrule
\end{tabularx}
\end{table}

The methodological lesson goes beyond these two fixes. A benchmark delimits what
it can see: measuring selection gave \SI{88}{\percent} success while execution of
the main task was failing in \SI{100}{\percent} of cases. \textbf{A flattering
indicator on one link says nothing about the chain.} The complete transcript of
the successful session is given in appendix~\ref{ann:session}.

\section{Evaluating the documentary assistant}
\label{sec:rag}

The documentary assistant was the last component of the system not supported by
any measurement. Its evaluation produced the sharpest result in this whole
chapter.

\subsection{Protocol}

What determines the quality of a documentary answer is the \textbf{retrieval}
that precedes it: if the wrong excerpt comes up, the answer will be wrong
whatever the quality of the language model. It is therefore retrieval alone that
is measured here.

Fifteen questions were written, each associated with the corpus document that
actually contains the answer. Two indicators are recorded: is the right document
the source of the first excerpt (accuracy@1), and does it appear among the four
retrieved excerpts (recall@4)? The second is the more relevant in practice,
since the agent receives all four excerpts.

As the corpus contains only four documents, a random draw would already give
about \SI{33}{\percent} accuracy@1. This threshold serves as a lower reference.

\subsection{A result below chance}

The first campaign gives \SI{27}{\percent} accuracy@1 — that is, \textbf{worse
than a random draw}. But the truly alarming indicator lies elsewhere: recall@4
is also \SI{27}{\percent}, exactly the same value.
